\section{Exact Formulations and Greedy Execution}
\label{sec:algorithms}

We separate three questions: how a deletion set is represented, how its
benefit is computed, and which optimization rule chooses the next tuple.
A compact representation does not by itself provide cheap marginal gains,
and an approximate marginal oracle need not preserve an exact-greedy
guarantee.

\subsection{A Unified Explicit-Derivation ILP}
\label{sec:ilp}

Introduce binary variables $x_t$ for $t\in\Del$, with $x_t=1$ denoting
deletion, and $z_b$ for $b\in\B$, with $z_b=1$ denoting survival.
Minimize
\begin{equation}
\label{eq:ilp-objective}
\sum_{t\in\Del}x_t
\end{equation}
subject to the linking constraints
\begin{align}
 z_b&\le 1-x_t &&(b\in\B,\ t\in T_b),\label{eq:ilp-upper}\\
 z_b&\ge 1-\sum_{t\in T_b}x_t &&(b\in\B).\label{eq:ilp-lower}
\end{align}
For \agg{COUNT} or \agg{SUM}, add
\begin{equation}
\label{eq:ilp-coverage}
\sum_{b\in\B}w_b(1-z_b)\ge\theta,
\end{equation}
where $w_b=1$ or $v_b$, respectively. For \agg{MAX} or \agg{MIN}, add
\begin{equation}
\label{eq:ilp-extrema}
z_b=0\quad(b\in\B_{\mathrm{crit}}),
\qquad \sum_{b\in\B}z_b\ge1.
\end{equation}
For downward \agg{AVG}, add
\begin{equation}
\label{eq:ilp-avg}
\sum_{b\in\B}(v_b-\tau_{\downarrow})z_b\le0,
\qquad \sum_{b\in\B}z_b\ge1.
\end{equation}
Binary rational coefficients may be cleared by positive denominators
without changing feasibility.

\begin{theorem}[Exactness of the ILP]
\label{thm:ilp}
The feasible ILP solutions correspond exactly to valid instability
witnesses through $S=\{t:x_t=1\}$. Its optimum is $|S^\star|$;
ILP infeasibility is equivalent to witness infeasibility.
\end{theorem}
\begin{proof}
If any $t\in T_b$ is deleted, constraint~\eqref{eq:ilp-upper} forces
$z_b=0$. If none is deleted, constraint~\eqref{eq:ilp-lower} forces
$z_b=1$. In particular, an empty support forces survival. Thus $z_b$
is precisely the actual survival indicator. Constraint~\eqref{eq:ilp-coverage}
then expresses the destroyed contribution. Lemma~\ref{lem:extrema-equivalence}
proves the extrema constraints, and multiplying the average inequality
by its positive denominator proves~\eqref{eq:ilp-avg}. Conversely,
assigning the true deletion and survival indicators of any valid witness
satisfies every constraint for its operator. The objective counts exactly
the deleted source tuples.
\end{proof}

Both linking directions matter. With values $\{10,0\}$, only the
high-valued row deletable, and target zero, the upper bounds alone would
allow setting both $x_{10}=0$ and $z_{10}=0$, inventing a free deletion.
The lower bound prevents this spurious solution.

The encoding has $n+M$ binary variables and $O(M+L)$ constraints.
It is polynomial in the explicit-derivation representation, not necessarily
in a supplied compact circuit. It provides an exact baseline when solved
to proven optimality. A timed-out feasible solution is an upper bound on
optimum; a solver's certified lower bound is needed to quantify its gap.
The use of deletion and witness variables follows established deletion
propagation ILP methodology~\cite{makhija2025ilp}; we do not claim the
existence of such an encoding as a new optimization technique, or transfer
integral-relaxation results to this formulation without proof.

\subsection{Explicit Supports and Incremental Coverage}
\label{sec:explicit-greedy}

For nonnegative coverage weights, store an incidence list of derivations
containing each tuple and a live flag for every derivation. Initialize
\[
g_t=\sum_{b:t\in T_b}w_b.
\]
At each iteration, choose an allowed undeleted tuple of maximum $g_t$.
For each still-live derivation containing it, mark that derivation dead
and subtract $w_b$ from the gain of every tuple in its support. Dead
derivations are never subtracted again. Maintain the accumulated destroyed
weight to test termination. For extrema, use only critical derivations,
unit weights, and the permitted ground set of the current protected run.

\begin{proposition}[Incremental marginal invariant]
\label{prop:incremental}
After any deletion set $S$ selected by this procedure, every remaining
allowed tuple satisfies $g_t=f(t\mid S)$. The procedure therefore
implements exact marginal-gain greedy.
\end{proposition}
\begin{proof}
Initially every derivation is live, so the formula is
Equation~\eqref{eq:coverage-marginal}. When a tuple is deleted, precisely
the previously live incident derivations cease to survive. Their weights
are removed once from every remaining tuple's gain, leaving exactly the
sum over still-surviving incident derivations.
\end{proof}

With a scan to select the next tuple, one run takes $O(M+L+n^2)$
arithmetic operations: the incidence lists are initialized once, each
derivation dies at most once, and its support is traversed once when it
dies. Scanning the chosen tuples' incidence lists costs at most $L$ over
the run. Space is $O(M+L+n)$. These costs exclude construction of the
derivation representation. The bound is for one protected run; gains and
live flags must be reinitialized for another survivor.

The disjoint-support method can be implemented by scanning reduced critical
supports while maintaining their selected union. Add a support exactly
when its intersection with that union is empty. One protected run takes
$O(M+L+n)$ time and space with indexed tuple membership. Hence the
bounded-support guarantee does not require solving an optimization
subproblem beyond the survivor enumeration.

\subsection{Query Oracles and Exact Lazy Selection}
\label{sec:query-oracles}

Let $C_Q$ bound an exact evaluation of the residual aggregate, including
arithmetic on its encoded values. For \agg{COUNT} and \agg{SUM},
\[
f(t\mid S)=Q(D\setminus S)-Q(D\setminus(S\cup\{t\})).
\]
Caching the current answer lets ordinary greedy compute all candidates by
at most $n$ residual evaluations per iteration. At most $n$ iterations
therefore require $O(n^2)$ evaluations and $O(n^2 C_Q)$ time. For extrema,
counting residual derivations satisfying the strict critical-value test
gives the corresponding $h$ oracle under bag semantics. Survivor enumeration
still requires a complete method of obtaining safe supports; value-oracle
access alone does not supply them.

Submodularity also permits the standard accelerated greedy
selection~\cite{minoux1978lazy}. Store a previously
computed marginal as an upper bound in a max-priority queue. Recompute
the largest bound at the current deletion set. Select its tuple only when
its exact gain is at least every other stored bound; otherwise update the
bound and continue. Since true gains can only decrease, that acceptance
condition certifies an actual maximum marginal. Fixing a common tuple
order resolves ties consistently with ordinary greedy. Refreshing each
tuple at most once per iteration gives the same $O(n^2)$ worst-case oracle
bound, with additional queue overhead; fewer evaluations are a possible
practical benefit, not a worst-case improvement established here.

Restricting candidates to an uncertified active set, using estimated gains,
or accepting a tuple before this maximum test defines a heuristic. Such
choices need a separate analysis and do not inherit
Theorem~\ref{thm:general-greedy} from the implementation alone.

\subsection{Arithmetic Provenance Circuits}
\label{sec:circuits}

A circuit oracle here is an acyclic arithmetic circuit with rational
constants, tuple-survival inputs $y_t$, and binary addition and
multiplication gates. Its size $c$ counts gates and wires. The oracle
contract is that, at every Boolean assignment, its output equals the
surviving weighted derivation sum:
\begin{equation}
\label{eq:circuit-contract}
P(y)=\sum_{b\in\B}w_b\prod_{t\in T_b}y_t
\qquad(y\in\{0,1\}^{n}).
\end{equation}
Here $y_t=0$ denotes deletion. The contract may be instantiated for SUM,
COUNT, or critical COUNT. Constructing a compact circuit satisfying it
is a separate task; an arbitrary circuit for query existence is
insufficient. This arithmetic representation is motivated by semiring
provenance~\cite{green2007provenance}.

A topological evaluation takes $O(c)$ arithmetic operations. At the
assignment $y^S$ encoding the current deletions, exact marginal gains are
\begin{equation}
\label{eq:circuit-difference}
f(t\mid S)=P(y^S)-P(y^S\text{ with }y_t=0).
\end{equation}
Thus any circuit meeting the Boolean contract implements exact greedy
in $O(n^2c)$ arithmetic operations per run, without expanding its
polynomial. Protected runs fix the survivor's inputs to one and exclude
them from deletion candidates. Arithmetic complexity excludes circuit
construction and must be augmented by operand bit lengths. The contract
alone does not bound those lengths for arbitrary supplied circuits.

\begin{proposition}[Exact derivatives for multilinear circuits]
\label{prop:gradient}
If the circuit's formal polynomial is multilinear in every deletable
variable, all current deletion marginals can be computed by one forward
and one reverse differentiation pass, using $O(c+n)$ arithmetic operations.
One greedy run then uses $O(n(c+n))$ such operations, including selection.
\end{proposition}
\begin{proof}
For any undeleted tuple $t$, multilinearity gives
$P(y)=y_t A(y_{-t})+B(y_{-t})$. The difference between setting $y_t$ to
one and to zero is $A(y_{-t})=\partial P/\partial y_t$. A forward pass
records gate values. The reverse pass uses standard arithmetic-circuit
differentiation~\cite{baur1983derivatives}. In reverse topological order, initialize the output's
adjoint to one; an addition gate passes its adjoint to both inputs, and a
product gate passes its adjoint times the other input's value. Sum
adjoints at shared gates and at all leaves with the same tuple label.
The chain rule gives all partial derivatives in one pass over gates and
wires, with no division and hence no difficulty at zero-valued inputs.
Scanning the at most $n$ candidates chooses an exact largest gain.
Repeating for at most $n$ iterations gives the stated cost.
\end{proof}

Formal multilinearity is stronger than Boolean agreement with
Equation~\eqref{eq:circuit-contract}. A provenance polynomial for a
self-join may contain $y_t^2$: its values at zero and one represent one
surviving derivation, but its derivative at one is two, whereas the
actual deletion loss is one. Ordinary differentiation therefore need
not compute exact gains for self-joins. Direct reevaluation via
Equation~\eqref{eq:circuit-difference} remains correct.

The usual derivation polynomial of a self-join-free query is multilinear
in tuple identifiers. A sufficient circuit-level certificate is that, at
every multiplication gate, the two child polynomials use disjoint sets of
deletable variables. Addition can share variables between its children.
This certificate can be checked by propagating variable sets through the
DAG, with a cost depending on their representation. We do not assume
that every factorized provenance representation has this property, nor
that converting one to a multilinear circuit preserves size.

\subsection{AVG Search and Verification}
\label{sec:avg-search}

At target $\tau=\tau_{\downarrow}$ define
\[
F_\tau(S)=\sum_{b\in\B(S)}(v_b-\tau).
\]
A valid AVG witness satisfies $F_\tau(S)\le0$ and $|\B(S)|\ge1$.
A natural heuristic chooses a tuple maximizing the net decrease
$F_\tau(S)-F_\tau(S\cup\{t\})$ among deletions leaving a survivor.
The gain must include all destroyed derivations, including those with
negative $v_b-\tau$; summing only positive coefficients is not the true
threshold gain. Signed weights invalidate the monotone-coverage argument.
A run may stall without proving infeasibility, so the ILP is the exact
fallback when a certificate is required. Independent-row instances instead
use Proposition~\ref{prop:avg-rows}.

Every returned candidate should be checked against the original residual
query with the prescribed bag semantics, deletability policy, and empty-input
condition. This check certifies feasibility, not optimality. In particular,
removing redundant deletions can yield an inclusion-minimal witness without
establishing minimum cardinality. Applying that postprocessing to a feasible
approximation can only improve its size, provided each removal preserves
validity and the threshold.
